\documentclass[12pt,a4paper]{article}

\usepackage[utf8]{inputenc}
\usepackage[T1]{fontenc}
\usepackage[english]{babel}
\usepackage{geometry}
\usepackage{graphicx}
\usepackage{booktabs}
\usepackage{longtable}
\usepackage{array}
\usepackage{hyperref}
\usepackage{listings}
\usepackage{xcolor}
\usepackage{enumitem}
\usepackage{amsmath}
\usepackage{float}

\geometry{margin=2.5cm}
\setlength{\parskip}{0.5em}
\setlength{\parindent}{0pt}

\hypersetup{
    colorlinks=true,
    linkcolor=blue,
    urlcolor=blue,
    citecolor=blue
}

\lstdefinestyle{json}{
    basicstyle=\ttfamily\footnotesize,
    breaklines=true,
    frame=single,
    backgroundcolor=\color{gray!8}
}
\lstdefinestyle{code}{
    basicstyle=\ttfamily\footnotesize,
    breaklines=true,
    frame=single,
    backgroundcolor=\color{gray!8},
    language=C++
}
\lstdefinestyle{bash}{
    basicstyle=\ttfamily\footnotesize,
    breaklines=true,
    frame=single,
    backgroundcolor=\color{gray!8},
    language=bash
}

\title{%
    \textbf{Quill}\\
    \large Post-kwantowy komunikator P2P z szyfrowaniem sesji\\
    \normalsize Dokumentacja techniczna projektu\\
    \small Studio Projektowe, semestr 6
}
\author{%
    AGH UST, Wydział Informatyki\\
    Kierunek: Informatyka i Systemy Inteligentne\\[0.8em]
    \textit{[Imię i nazwisko autora/autorów]}
}
\date{Czerwiec 2026 \quad wersja 1.0.0}

\begin{document}

\maketitle

\begin{abstract}
\noindent
\textbf{Quill} to komunikator sieciowy napisany w C++23, w którym cała warstwa asymetryczna opiera się na standardach post-kwantowych NIST (FIPS~203 ML-KEM, FIPS~204 ML-DSA). Projekt implementuje własny protokół handshake, szyfrowanie sesji AES-256-GCM z ochroną przed replay, mechanizm TOFU, trwałe profile użytkownika oraz transfer plików z integralnością SHA-3. Aplikacja posiada interfejs ImGui oraz bibliotekę rdzeniową \texttt{quill\_core} testowaną 97 testami Google Test i weryfikowaną w CI na GitHub Actions.

\medskip
\noindent\textbf{Słowa kluczowe:} kryptografia post-kwantowa, ML-KEM, ML-DSA, AES-GCM, TOFU, komunikator P2P, C++23.
\end{abstract}

\tableofcontents
\newpage

% ═══════════════════════════════════════════════════════════════════
\section{Wprowadzenie}
% ═══════════════════════════════════════════════════════════════════

\subsection{Kontekst projektu}

Projekt \textbf{Quill} powstał w ramach przedmiotu \emph{Studio Projektowe} na 6. semestrze studiów inżynierskich na Akademii Górniczo-Hutniczej w Krakowie, na kierunku Informatyka i Systemy Inteligentne. Celem było zaprojektowanie i implementacja działającego systemu komunikacji, który nie opiera się na klasycznej kryptografii asymetrycznej podatnej na ataki kwantowe, lecz na algorytmach zatwierdzonych przez NIST w 2024 roku.

Quill nie jest wrapperem TLS ani gotowym protokołem aplikacyjnym --- implementuje własny stos od warstwy transportowej (TCP) po szyfrowanie wiadomości, z pełną kontrolą nad każdym krokiem handshake i każdym bajtem na linii. Taki poziom szczegółowości jest celowy: projekt ma być obronny akademicko i stanowić fundament planowanej pracy inżynierskiej na temat inteligentnego doboru algorytmów PQC w zależności od kontekstu (urządzenie, opóźnienie, wymagana siła kryptograficzna).

\subsection{Cel i zakres}

\textbf{Cel główny:} zbudować komunikator, w którym:
\begin{enumerate}[label=\alph*)]
    \item wymiana klucza i uwierzytelnianie opierają się wyłącznie na algorytmach post-kwantowych,
    \item sesja jest szyfrowana AES-256-GCM z wiązaniem kontekstu (AAD) i ochroną przed replay,
    \item użytkownik ma trwałą tożsamość kryptograficzną chronioną hasłem na dysku,
    \item zaufanie do peera jest zarządzane mechanizmem TOFU (jak w SSH),
    \item możliwy jest bezpieczny transfer plików z weryfikacją integralności.
\end{enumerate}

\textbf{Zakres zrealizowany} obejmuje m.in.: handshake PQC, trzy poziomy bezpieczeństwa (FAST/BALANCED/MAX), Perfect Forward Secrecy, pokoje rozmów, interfejs graficzny z narzędziami audytowymi, 97 testów automatycznych oraz pipeline CI.

\textbf{Poza zakresem} (świadomie odłożone): NAT traversal (UDP hole punching), prawdziwy end-to-end bez zaufania do serwera-relay, wersjonowanie schematu JSON protokołu, streaming plików bez buforowania w RAM.

\subsection{Struktura dokumentu}

Dokument opisuje kolejno: tło teoretyczne i model zagrożeń (rozdz.~2--3), architekturę systemu (rozdz.~4), stos kryptograficzny (rozdz.~5), szczegółowy protokół komunikacji (rozdz.~6), warstwę sieciową (rozdz.~7), tożsamość i zaufanie (rozdz.~8), moduły implementacyjne (rozdz.~9), interfejs użytkownika (rozdz.~10), scenariusze użycia (rozdz.~11), testowanie (rozdz.~12), budowę i wdrożenie (rozdz.~13), analizę bezpieczeństwa (rozdz.~14), ograniczenia (rozdz.~15), podsumowanie (rozdz.~16) oraz załączniki.

% ═══════════════════════════════════════════════════════════════════
\section{Motywacja i tło teoretyczne}
% ═══════════════════════════════════════════════════════════════════

\subsection{Zagrożenie kwantowe}

Klasyczna kryptografia asymetryczna (RSA, DSA, ECDH, ECDSA) opiera się na problemach, które na komputerze kwantowym rozwiązuje algorytm Shora w czasie wielomianowym. Oznacza to, że długoterminowo przechwycone dzisiejsze sesje zaszyfrowane RSA/ECDH mogą zostać odszyfrowane po pojawieniu się wystarczająco dużego komputera kwantowego (\emph{harvest now, decrypt later}).

W 2024 roku NIST opublikował pierwsze standardy kryptografii post-kwantowej:
\begin{itemize}
    \item \textbf{FIPS 203} --- ML-KEM (Module-Lattice Key Encapsulation Mechanism), wcześniej znany jako CRYSTALS-Kyber,
    \item \textbf{FIPS 204} --- ML-DSA (Module-Lattice Digital Signature Algorithm), wcześniej CRYSTALS-Dilithium,
    \item \textbf{FIPS 205} --- SLH-DSA (SPHINCS+) --- podpisy hash-based (w Quill nieużywane bezpośrednio, ale dostępne w liboqs).
\end{itemize}

Quill implementuje ML-KEM i ML-DSA (oraz Falcon-512 w poziomie FAST) jako zamienniki całej klasycznej asymetrii.

\subsection{Algorytm Grovera a AES}

Symetryczne szyfrowanie AES-256 nie jest podatne na Shora w taki sam sposób jak RSA. Algorytm Grovera redukuje efektywną siłę klucza o połowę --- AES-256 daje $\sim$128 bitów bezpieczeństwa kwantowego, co nadal jest uznawane za bezpieczne. Dlatego Quill zachowuje AES-256-GCM jako szyfr sesji, koncentrując wysiłek PQC na wymianie klucza i podpisach.

\subsection{Porównanie stosów kryptograficznych}

\begin{table}[H]
\centering
\begin{tabular}{@{}p{3.5cm}p{4.5cm}p{5.5cm}@{}}
\toprule
\textbf{Funkcja} & \textbf{Klasycznie} & \textbf{Quill (PQC)} \\
\midrule
Wymiana klucza & ECDH, RSA-OAEP & ML-KEM (Kyber-512/768/1024) \\
Podpis / auth & ECDSA, RSA-PSS & ML-DSA-65/87, Falcon-512 \\
Szyfrowanie sesji & AES-256-GCM & AES-256-GCM (bez zmian) \\
KDF & często HKDF / TLS PRF & HKDF-SHA256 (RFC 5869) \\
Hash integralności & SHA-256 & SHA-3-256 \\
\bottomrule
\end{tabular}
\caption{Mapowanie funkcji kryptograficznych w Quill.}
\end{table}

\subsection{Podstawy kryptografii na kratach}

Kyber (ML-KEM) i Dilithium (ML-DSA) opierają się na trudności problemów na strukturalnych kratach (Module-LWE, Module-SIS). Intuicyjnie: klucz publiczny to „zaszumione'' równania liniowe, a klucz prywatny pozwala odszumiać wynik. Bezpieczeństwo opiera się na założeniu, że najlepsze znane ataki wymagają wykładniczego czasu na klasycznym i kwantowym komputerze (w odróżnieniu od faktoryzacji i logarytmu dyskretnego).

Quill nie implementuje tych algorytmów od zera --- używa biblioteki \textbf{liboqs} (Open Quantum Safe), która dostarcza referencyjne implementacje zweryfikowane w konkursie NIST. Własna implementacja obejmuje natomiast protokół, KDF, TOFU, serializację i logikę aplikacji.

% ═══════════════════════════════════════════════════════════════════
\section{Model zagrożeń i założenia}
% ═══════════════════════════════════════════════════════════════════

\subsection{Zagrożenia uwzględnione}

\begin{table}[H]
\centering
\small
\begin{tabular}{@{}p{3.5cm}p{9cm}@{}}
\toprule
\textbf{Zagrożenie} & \textbf{Mechanizm obrony w Quill} \\
\midrule
Podsłuch na łączu & AES-256-GCM na każdej wiadomości; klucz sesji z ML-KEM \\
Atak MITM przy handshake & Podpisy ML-DSA/Falcon + TOFU fail-closed \\
Replay wiadomości czatu & Monotoniczny seq wiązany przez GCM AAD \\
Manipulacja ciphertextu & Tag GCM 128-bit; odrzucenie przy weryfikacji \\
Zmiana klucza peera & TOFU MISMATCH --- blokada przed KEM encaps \\
Kradzież pliku klucza z dysku & Argon2id + AES-256-GCM (format QID2) \\
Uszkodzenie danych w transferze pliku & SHA-3 per chunk + hash całości w FILE\_END \\
Wstrzyknięcie chunka z innego transferu & AAD \texttt{FILE|tid|index} + chunk\_hash \\
\bottomrule
\end{tabular}
\caption{Zagrożenia i odpowiadające im mechanizmy.}
\end{table}

\subsection{Założenia zaufania}

\begin{itemize}
    \item Użytkownik chroni passphrase profilu --- to jedyna obrona offline przy kradzieży pliku klucza.
    \item Porównanie fingerprintów out-of-band (telefon, QR) jest odpowiedzialnością użytkownika --- TOFU sam nie gwarantuje weryfikacji tożsamości przy pierwszym kontakcie.
    \item \textbf{Serwer jest zaufanym relayem} --- widzi plaintext wiadomości. Każdy hop klient$\leftrightarrow$serwer jest szyfrowany, ale serwer odszyfrowuje i ponownie szyfruje dla odbiorców.
    \item Sieć TCP jest niezawodna w dostarczaniu bajtów w kolejności (brak UDP, brak NAT traversal).
\end{itemize}

\subsection{Co Quill \emph{nie} chroni}

\begin{itemize}
    \item Atak na działający proces (keyloggery, pamięć RAM po odblokowaniu profilu).
    \item Kompromitacja serwera-hosta (serwer widzi treść wiadomości).
    \item Denial of service (brak rate limiting, brak ochrony przed floodem TCP).
    \item Ukrywanie metadanych (adresy IP, czasy połączeń, rozmiary pakietów).
\end{itemize}

% ═══════════════════════════════════════════════════════════════════
\section{Architektura systemu}
% ═══════════════════════════════════════════════════════════════════

\subsection{Przegląd warstw}

System podzielony jest na warstwę prezentacji (ImGui), logikę aplikacji (\texttt{ChatApp}) oraz bibliotekę rdzeniową \texttt{quill\_core}, która nie zależy od GUI.

\begin{lstlisting}[style=code,caption={Hierarchia modułów.}]
main.cpp
  └── ChatApp (frontend/)
        ├── ProfileManager, TrustStore (UI + aktywacja)
        ├── NetworkServer / NetworkClient (wątki sieciowe)
        └── FileReceiver, wizualizacja handshake
              │
              ▼
        quill_core (biblioteka statyczna)
              ├── crypto/    CryptoManager, KyberKEM, DilithiumSign,
              │              Hkdf, AesGcm, Argon2, IdentityManager
              ├── network/   Socket, framing TCP
              └── protocol/  MessageFormat, FileTransfer
\end{lstlisting}

\textbf{Decyzja projektowa:} separacja \texttt{quill\_core} od UI umożliwia uruchomienie 97 testów i CI bez GLFW/OpenGL/ImGui. Ta sama logika kryptograficzna działa w aplikacji i w testach integracyjnych (\texttt{socketpair}).

\subsection{Model hub-and-spoke}

Quill nie jest czystym P2P mesh --- działa w modelu \emph{gwiazdy}:
\begin{itemize}
    \item jedna instancja uruchamia \textbf{serwer} (nasłuch TCP),
    \item pozostałe instancje łączą się jako \textbf{klienci},
    \item serwer przekazuje wiadomości między użytkownikami w pokojach (rooms),
    \item każdy klient ma osobny klucz sesji AES z serwerem (osobny handshake).
\end{itemize}

Zalety: prosty broadcast, jeden punkt spotkania, łatwe demo na jednej maszynie (localhost). Wada: serwer widzi plaintext --- opisana w rozdz.~15.

\subsection{Wielowątkowość}

Serwer obsługuje wielu klientów --- każde połączenie TCP ma dedykowany wątek \texttt{std::thread} (\texttt{server\_client\_handler}). Główny wątek renderuje ImGui (60 FPS). Współdzielony stan (lista klientów, pokoje) chroniony jest mutexem \texttt{m\_clients\_mtx}. Licznik wiadomości \texttt{m\_msg\_count} jest atomowy.

\subsection{Pokoje (rooms) i izolacja wiadomości}

Serwer utrzymuje mapę \texttt{m\_rooms}: nazwa pokoju $\rightarrow$ zbiór klientów. Domyślny pokój to \texttt{general}. Klient wybiera pokój w UI; serwer przypisuje go w strukturze \texttt{ConnectedClient::room}.

\textbf{Broadcast} (\texttt{broadcast\_to\_room}) wysyła wiadomość do wszystkich uczestników pokoju oprócz nadawcy. Wiadomości z innych pokoi nie docierają do odbiorców --- izolacja na poziomie aplikacji (serwer filtruje po polu \texttt{room}).

Każdy klient w pokoju ma osobny klucz sesji AES (osobny handshake przy połączeniu). Serwer musi więc odszyfrować wiadomość od nadawcy i zaszyfrować ją osobno dla każdego odbiorcy --- stąd model hub-and-spoke opisany w rozdz.~4.

\subsection{Historia decyzji projektowych}

\begin{table}[H]
\centering
\small
\begin{tabular}{@{}p{4cm}p{4cm}p{5.5cm}@{}}
\toprule
\textbf{Decyzja} & \textbf{Alternatywa} & \textbf{Uzasadnienie} \\
\midrule
Własny protokół & TLS 1.3 + PQC hybrid & Pełna kontrola, cel edukacyjny \\
liboqs zamiast własnego KEM & Implementacja Kyber od zera & Bezpieczeństwo, czas \\
TOFU zamiast PKI & Certyfikaty X.509 & Zakres semestralny, jak SSH \\
JSON na TCP & Protobuf / binarny & Czytelność, debug, Packet Inspector \\
ImGui zamiast Qt & Qt Widgets / web UI & Lekkość, szybki prototyp \\
quill\_core bez UI & Monolit & Testowalność, CI headless \\
Hub relay & Mesh P2P & Prostszy broadcast, demo 1-serwer \\
\bottomrule
\end{tabular}
\caption{Wybrane decyzje architektoniczne.}
\end{table}

\subsection{Struktura katalogów projektu}

\begin{lstlisting}[style=bash,caption={Drzewo repozytorium (skrót).}]
Quill/
  main.cpp                 # entry point GLFW + ImGui
  CMakeLists.txt
  src/
    crypto/                # 10 modułów kryptograficznych
    network/               # Socket, Server, Client
    protocol/              # MessageFormat, FileTransfer
    frontend/              # ChatApp, Theme, render
  tests/                   # 11 plików testowych, 97 przypadków
  third_party/
    imgui/                 # git submodule v1.92.8
    portable-file-dialogs/   # wybór plików w GUI
  docs/                    # dokumentacja (ten plik)
  .github/workflows/ci.yml
\end{lstlisting}

% ═══════════════════════════════════════════════════════════════════
\section{Stos kryptograficzny}
% ═══════════════════════════════════════════════════════════════════

\subsection{Biblioteki zewnętrzne}

\begin{table}[H]
\centering
\small
\begin{tabular}{@{}p{2.8cm}p{4cm}p{6cm}@{}}
\toprule
\textbf{Biblioteka} & \textbf{Licencja} & \textbf{Rola w Quill} \\
\midrule
liboqs & MIT & ML-KEM, ML-DSA, Falcon \\
OpenSSL & Apache 2.0 & AES-GCM, HKDF, SHA-3, RAND \\
libargon2 & CC0/Apache & Fallback Argon2id (OpenSSL $<$ 3.2) \\
nlohmann/json & MIT & Serializacja wiadomości \\
Dear ImGui & MIT & Interfejs graficzny \\
GLFW + OpenGL & zlib / --- & Okno i kontekst renderowania \\
Google Test & BSD-3 & Framework testów \\
\bottomrule
\end{tabular}
\end{table}

\subsection{Poziomy bezpieczeństwa}

Enum \texttt{SecurityLevel} definiuje trzy tryby używane przy połączeniu:

\begin{table}[H]
\centering
\begin{tabular}{@{}llllp{3.5cm}@{}}
\toprule
\textbf{Poziom} & \textbf{KEM} & \textbf{Podpis} & \textbf{NIST} & \textbf{Zastosowanie} \\
\midrule
\texttt{FAST} & Kyber-512 & Falcon-512 & Level 1 & IoT, niskie opóźnienie \\
\texttt{BALANCED} & Kyber-768 & ML-DSA-65 & Level 3 & Domyślny, kompromis \\
\texttt{MAX} & Kyber-1024 & ML-DSA-87 & Level 5 & Dane krytyczne \\
\bottomrule
\end{tabular}
\caption{Mapowanie poziomów na algorytmy liboqs.}
\end{table}

Poziom wybierany jest w UI przed połączeniem i musi być zgodny po obu stronach (negocjacja w handshake). Benchmarki w aplikacji mierzą czasy keygen/encaps/decaps/sign/verify dla każdego poziomu.

\subsection{KyberKEM --- wymiana klucza}

Klasa \texttt{KyberKEM} opakowuje liboqs:
\begin{itemize}
    \item \texttt{generate\_keypair()} --- generacja efemerycznej pary na handshake (PFS),
    \item \texttt{encapsulate(pub)} $\rightarrow$ (ciphertext, shared\_secret),
    \item \texttt{decapsulate(ct, sk)} $\rightarrow$ shared\_secret.
\end{itemize}

Nowa para Kyber przy każdym handshake oznacza, że kompromitacja klucza sesji nie ujawnia przyszłych sesji (\emph{forward secrecy} na poziomie KEM). Klucz DSA pozostaje trwały --- służy tożsamości, nie sekretowi sesji.

\subsection{DilithiumSign --- podpisy}

Klasa \texttt{DilithiumSign} obsługuje ML-DSA-65/87 oraz Falcon-512 (poziom FAST):
\begin{itemize}
    \item \texttt{sign(message, secret\_key)} $\rightarrow$ podpis,
    \item \texttt{verify(message, signature, public\_key)} $\rightarrow$ bool,
    \item weryfikacja \textbf{zawsze} przed użyciem materiału KEM (inwariant w \texttt{CryptoManager}).
\end{itemize}

Podpisy wiążą: (1) klucz publiczny Kyber serwera w SRV\_HELLO, (2) ciphertext KEM klienta w CLI\_KEX.

\subsection{Hkdf --- wyprowadzanie klucza sesji}

Surowy \texttt{shared\_secret} z ML-KEM nie jest używany bezpośrednio jako klucz AES. HKDF-SHA256 (RFC 5869, OpenSSL EVP\_KDF) wyprowadza 32-bajtowy klucz:

\begin{itemize}
    \item \textbf{IKM:} shared\_secret z Kyber,
    \item \textbf{salt:} transkrypt handshake = \texttt{kyber\_pub $\Vert$ kem\_ciphertext},
    \item \textbf{info:} \texttt{quill-aes256gcm-session-v1|<LEVEL>} --- separacja domen,
    \item po HKDF: \texttt{OPENSSL\_cleanse} na shared\_secret.
\end{itemize}

Wiązanie z transkryptem gwarantuje, że klucz sesji jest unikalny dla \emph{tej konkretnej} wymiany, nawet gdyby ten sam shared\_secret powtórzył się (teoretycznie nieprawdopodobne, ale wymóg defensywny).

\subsection{AesGcm --- szyfrowanie sesji}

AES-256-GCM przez OpenSSL EVP:
\begin{itemize}
    \item nonce 12 bajtów losowych (\texttt{RAND\_bytes}) na każde szyfrowanie,
    \item tag uwierzytelniający 16 bajtów,
    \item opcjonalny AAD (Additional Authenticated Data) --- używany dla CHAT i FILE\_CHUNK,
    \item odrzucenie przy złym kluczu, zmanipulowanym ciphertext/tag lub niespójnym AAD.
\end{itemize}

\subsection{Argon2 --- ochrona kluczy na dysku}

Format \textbf{QID2} szyfruje klucz prywatny DSA:
\begin{itemize}
    \item Argon2id: $t=3$ iteracje, $m=64$\,MiB pamięci (RFC 9106),
    \item losowa sól 16 B, parametry zapisane w pliku,
    \item wyprowadzony klucz 32 B $\rightarrow$ AES-256-GCM,
    \item złe hasło = błąd weryfikacji tagu GCM (jeden komunikat dla złego hasła i manipulacji).
\end{itemize}

Na systemach z OpenSSL $\geq$ 3.2 używany jest \texttt{EVP\_KDF} ARGON2ID; na starszych (Ubuntu w CI) --- biblioteka \texttt{libargon2}.

% ═══════════════════════════════════════════════════════════════════
\section{Protokół komunikacji}
% ═══════════════════════════════════════════════════════════════════

\subsection{Format na linii (wire format)}

Każda wiadomość na TCP ma prefiks długości:
\begin{enumerate}
    \item 4 bajty \texttt{uint32\_t} --- długość payloadu (little-endian),
    \item $N$ bajtów --- payload (JSON tekstowy lub binarny zakodowany w JSON).
\end{enumerate}

Funkcje \texttt{Socket::send\_bytes} / \texttt{receive\_bytes} realizują semantykę \emph{send all} / \emph{recv all} --- gwarantują dostarczenie kompletnej ramki.

\subsection{Typy wiadomości}

Enum \texttt{MsgType} w \texttt{MessageFormat.h}:

\begin{table}[H]
\centering
\begin{tabular}{@{}llp{7.5cm}@{}}
\toprule
\textbf{Typ} & \textbf{Faza} & \textbf{Opis} \\
\midrule
\texttt{HANDSHAKE} & Negocjacja & SRV\_HELLO, CLI\_KEX (JSON bez szyfrowania) \\
\texttt{CHAT} & Sesja & Zaszyfrowana wiadomość tekstowa \\
\texttt{FILE\_START} & Transfer & Metadane pliku \\
\texttt{FILE\_CHUNK} & Transfer & Zaszyfrowany fragment 64\,KB \\
\texttt{FILE\_END} & Transfer & Hash SHA-3-256 całego pliku \\
\texttt{FILE\_NACK} & Transfer & Lista brakujących chunk\_index \\
\texttt{SYSTEM\_CMD} & Sterowanie & Polecenia wewnętrzne (PFS, pokoje) \\
\bottomrule
\end{tabular}
\end{table}

Struktura \texttt{SecureMessage} serializuje pola do JSON. Dane binarne (nonce, payload, hash) kodowane są jako tablice bajtów --- \emph{nie} jako \texttt{json::binary}, które nie przeżywa \texttt{dump()}/\texttt{parse()} (bug naprawiony i pokryty testami).

\subsection{Handshake PQC --- przebieg szczegółowy}

\subsubsection{Diagram sekwencji}

\begin{lstlisting}[caption={Handshake Quill (uproszczony).}]
KLIENT (Alice)                              SERWER (Bob)
      |                                           |
      |<-------- SRV_HELLO -----------------------|
      |   sig_pub_B, kyber_pub_B, sig_B(kyber)    |
      |                                           |
      | verify sig_B                              |
      | TOFU check (fail-closed)                  |
      | KEM encaps(kyber_pub_B) -> ct, ss         |
      | sign_A(ct)                                |
      |-------- CLI_KEX ------------------------->|
      |   sig_pub_A, ciphertext, sig_A(ct)        |
      |                                           | verify sig_A
      |                                           | TOFU remember
      |                                           | KEM decaps -> ss
      |                                           |
      | obie strony: key = HKDF(ss, salt, info)   |
      |======== sesja AES-256-GCM ================|
\end{lstlisting}

\subsubsection{SRV\_HELLO (serwer $\rightarrow$ klient)}

Przykładowa struktura JSON (pola binarne skrócone):

\begin{lstlisting}[style=json]
{
  "type": "SRV_HELLO",
  "sig_pub":   [/* bajty klucza publicznego DSA */],
  "kyber_pub": [/* bajty klucza publicznego Kyber */],
  "sig":       [/* podpis kyber_pub kluczem DSA serwera */]
}
\end{lstlisting}

\subsubsection{CLI\_KEX (klient $\rightarrow$ serwer)}

\begin{lstlisting}[style=json]
{
  "type": "CLI_KEX",
  "sig_pub":   [/* klucz publiczny DSA klienta */],
  "ciphertext":[/* wynik KEM encaps */],
  "sig":       [/* podpis ciphertext kluczem DSA klienta */]
}
\end{lstlisting}

\subsubsection{Inwarianty bezpieczeństwa handshake}

\begin{enumerate}
    \item \texttt{verify()} podpisu \textbf{przed} \texttt{encaps()} / \texttt{decaps()}.
    \item TOFU \texttt{MISMATCH} $\rightarrow$ wyjątek \texttt{TofuMismatchError} \textbf{przed} wysłaniem ciphertextu KEM.
    \item Klucz sesji wyłącznie przez HKDF; shared\_secret czyszczony z pamięci.
    \item Kyber keypair efemeryczny; DSA z \texttt{IdentityManager::load\_or\_generate}.
\end{enumerate}

\subsection{Wiadomości CHAT}

Po handshake każda wiadomość czatu:

\begin{enumerate}
    \item Nadawca inkrementuje \texttt{send\_seq}, buduje AAD = \texttt{"CHAT|" + std::to\_string(seq)}.
    \item Plaintext szyfrowany AES-256-GCM z losowym nonce 12 B.
    \item JSON zawiera: \texttt{type}, \texttt{sender}, \texttt{seq}, \texttt{nonce}, \texttt{payload} (ciphertext+tag), \texttt{timestamp}.
    \item Odbiorca deszyfruje, weryfikuje AAD, odrzuca jeśli \texttt{seq <= recv\_seq}.
\end{enumerate}

Ochrona przed replay jest \emph{fail-closed}: podmiana numeru sekwencji w JSON bez ponownego szyfrowania psuje tag GCM (AAD musi się zgadzać). Ponowne wysłanie starej wiadomości z tym samym seq jest odrzucane logicznie.

\textbf{Reset liczników:} przy każdym nowym handshake i rotacji PFS liczniki \texttt{send\_seq}/\texttt{recv\_seq} zerowane --- stary szyfrogram i tak nie deszyfruje się nowym kluczem.

\subsection{Perfect Forward Secrecy (PFS)}

Rotacja kluczy inicjowana jest:
\begin{itemize}
    \item ręcznie z UI (przycisk Rotate Keys),
    \item automatycznie na serwerze co $N$ wiadomości (konfigurowalne).
\end{itemize}

Przebieg: nowy handshake PQC na istniejącym połączeniu TCP (bez disconnect). Serwer ustawia flagę \texttt{pending\_pfs\_rotation}; wątek handlera klienta wykonuje \texttt{perform\_pfs\_rotation}. \texttt{Socket::try\_receive\_bytes(500ms)} pozwala obsłużyć rotację bez blokowania na brak wiadomości od klienta.

\subsection{Transfer plików}

\subsubsection{Przebieg protokołu}

\begin{enumerate}
    \item \textbf{FILE\_START} --- \texttt{transfer\_id} (UUID), \texttt{file\_name}, \texttt{file\_size}, \texttt{total\_chunks}.
    \item \textbf{FILE\_CHUNK} $\times N$ --- każdy chunk max 64\,KB plaintextu:
    \begin{itemize}
        \item szyfrowanie AES-256-GCM z unikalnym nonce,
        \item AAD = \texttt{"FILE|" + transfer\_id + "|" + chunk\_index},
        \item \texttt{chunk\_hash} = SHA3-256(plaintext chunka) w JSON --- weryfikacja \emph{przed} buforowaniem,
    \end{itemize}
    \item \textbf{FILE\_END} --- \texttt{file\_hash} = SHA3-256(cały plik).
    \item \textbf{FILE\_NACK} (opcjonalnie) --- \texttt{missing: [indeksy]}; nadawca retransmituje z \textbf{nowym nonce} (max 5 rund).
\end{enumerate}

\subsubsection{Selective repeat}

Klasa \texttt{FileSenderSession} trzyma plaintext wszystkich chunków w pamięci do momentu zakończenia transferu (umożliwia retransmisję). \texttt{FileReceiver} buforuje chunki w mapie \texttt{chunk\_index $\rightarrow$ bytes} i scala przez \texttt{assemble()} po kompletności.

Jeśli \texttt{FILE\_END} dotrze przed wszystkimi chunkami, odbiorca ustawia \texttt{awaiting\_end = true} i czeka na uzupełnienie przez NACK. Po przekroczeniu 5 rund NACK transfer kończy się błędem.

\subsubsection{Relay serwera}

Serwer odszyfrowuje chunk od nadawcy i szyfruje ponownie dla każdego odbiorcy w pokoju tym samym kluczem sesji co czat, zachowując ten sam AAD i \texttt{chunk\_hash}.

% ═══════════════════════════════════════════════════════════════════
\section{Warstwa sieciowa}
% ═══════════════════════════════════════════════════════════════════

\subsection{Klasa Socket}

Abstrakcja nad deskryptorem POSIX:
\begin{itemize}
    \item \texttt{send\_bytes} / \texttt{receive\_bytes} --- blokujące I/O z framingiem,
    \item \texttt{try\_receive\_bytes(timeout\_ms)} --- nieblokujące z timeoutem (PFS, cleanup),
    \item RAII --- zamknięcie FD w destruktorze.
\end{itemize}

\subsection{NetworkServer}

\begin{itemize}
    \item \texttt{bind(host, port)} + \texttt{listen()},
    \item pętla \texttt{accept()} w wątku głównym serwera,
    \item każdy klient: nowy \texttt{Socket}, handshake przez \texttt{CryptoManager::server\_handshake}, rejestracja w \texttt{m\_clients},
    \item \texttt{broadcast\_to\_room(room, packet)} --- wysyłka do wszystkich w pokoju oprócz nadawcy,
    \item cleanup martwych połączeń: usunięcie z \texttt{m\_clients} i \texttt{m\_rooms}.
\end{itemize}

\subsection{NetworkClient}

\begin{itemize}
    \item \texttt{connect(host, port)},
    \item handshake przez \texttt{CryptoManager::client\_handshake} z \texttt{peer\_id = "srv:" + host + ":" + port},
    \item pętla odbioru w wątku roboczym, kolejka wiadomości do UI.
\end{itemize}

\subsection{Ograniczenia sieciowe}

Quill używa wyłącznie TCP IPv4/IPv6. Brak:
\begin{itemize}
    \item UDP hole punching,
    \item serwera rendezvous,
    \item TLS wrappera (niepotrzebny --- własny PQC handshake),
    \item kompresji payloadu.
\end{itemize}

Oba peeri muszą być osiągalne pod znanym adresem IP i portem (np.\ localhost dla demo, sieć LAN dla prezentacji).

% ═══════════════════════════════════════════════════════════════════
\section{Tożsamość, profile i zaufanie}
% ═══════════════════════════════════════════════════════════════════

\subsection{IdentityManager}

Trwała tożsamość kryptograficzna --- para kluczy DSA generowana raz i przechowywana na dysku.

\textbf{Lokalizacja:} \texttt{\textasciitilde/.quill/profiles/<nazwa>/} (aktywowana przez profil) lub fallback \texttt{\$QUILL\_IDENTITY\_DIR} / \texttt{\textasciitilde/.quill/identity}.

\textbf{Bezpieczeństwo plików:}
\begin{itemize}
    \item katalog: chmod 0700,
    \item plik klucza: chmod 0600,
    \item zapis atomowy: plik tymczasowy + \texttt{rename()},
    \item odmowa wczytania przy zbyt szerokich uprawnieniach (wzorzec OpenSSH),
    \item self-test sign/verify po wczytaniu --- uszkodzony klucz = wyjątek, nie cicha regeneracja,
    \item suma SHA-3-256 payloadu w pliku QID1 --- detekcja korupcji bajtów.
\end{itemize}

\textbf{Formaty:}
\begin{itemize}
    \item \textbf{QID1} --- klucz plaintext (legacy, ostrzeżenie w UI),
    \item \textbf{QID2} --- Argon2id + AES-256-GCM; automatyczna migracja QID1$\rightarrow$QID2 przy odblokowaniu hasłem.
\end{itemize}

\textbf{Fingerprint:} SHA-3-256 klucza publicznego, pierwsze 12 bajtów w formacie \texttt{XXXX-XXXX-XXXX-XXXX-XXXX-XXXX} (np.\ \texttt{A3F2-9C1B-44DE-F021-8B3C-7A09}).

\subsection{ProfileManager}

Profil = katalog z \texttt{profile.json} + pliki kluczy.

\textbf{Tworzenie profilu:}
\begin{itemize}
    \item nazwa: whitelist znaków (ochrona przed path traversal),
    \item passphrase: min.\ 8 znaków, odrzucenie \texttt{aaaaaaaa} itp.,
    \item nowe profile zawsze \texttt{encrypted=true} (QID2).
\end{itemize}

\textbf{Zmienne środowiskowe:}
\begin{itemize}
    \item \texttt{QUILL\_PROFILES\_DIR} --- nadpisanie \texttt{\textasciitilde/.quill/profiles},
    \item \texttt{QUILL\_IDENTITY\_DIR} --- katalog tożsamości (fallback).
\end{itemize}

\textbf{Logout:} czyści passphrase i klucze tajne z pamięci (\texttt{OPENSSL\_cleanse}); zablokowany w trakcie aktywnej sesji sieciowej (klucze potrzebne do PFS).

\subsection{TrustStore (TOFU)}

Plik \texttt{known\_hosts} w katalogu profilu (JSON, chmod 0600).

\textbf{Identyfikacja peerów:}
\begin{itemize}
    \item klient $\rightarrow$ serwer: \texttt{"srv:<host>:<port>"},
    \item serwer $\rightarrow$ klient: \texttt{"cli:<fingerprint>"} (klient nie ma stałego adresu).
\end{itemize}

\textbf{Stany:}
\begin{description}[style=nextline]
    \item[UNVERIFIED] Pierwszy kontakt; klucz zapisany; ostrzeżenie w UI.
    \item[KNOWN] Klucz zgodny z zapisem; fingerprint nie potwierdzony out-of-band.
    \item[VERIFIED] Użytkownik kliknął „Mark as Verified'' po porównaniu fingerprintu.
    \item[MISMATCH] Inny klucz niż zapisany --- połączenie zablokowane; wpis \emph{nie} nadpisywany.
\end{description}

\textbf{Usunięcie wpisu} (\texttt{TrustStore::remove}) --- jedyna droga po MISMATCH do ponownego połączenia ze zmienionym kluczem (świadoma decyzja użytkownika).

Porównanie kluczy odbywa się na pełnym SHA-3-256, nie na skróconym fingerprintcie wyświetlanym w UI.

% ═══════════════════════════════════════════════════════════════════
\section{Moduły implementacyjne --- szczegóły}
% ═══════════════════════════════════════════════════════════════════

\subsection{CryptoManager}

Centralny moduł handshake. Metody publiczne:
\begin{itemize}
    \item \texttt{client\_handshake(sock, peer\_id, on\_step)},
    \item \texttt{server\_handshake(sock, on\_step)}.
\end{itemize}

Zwraca \texttt{HandshakeResult}: \texttt{session\_key} (32 B), \texttt{peer\_fingerprint}, \texttt{peer\_trust}, \texttt{total\_ms}.

\texttt{ChatApp::do\_client\_handshake} i \texttt{do\_server\_handshake} to cienkie wrappery dodające logi i wizualizator kroków --- cała logika bezpieczeństwa jest w \texttt{quill\_core}.

\subsection{MessageFormat}

Serializacja/deserializacja \texttt{SecureMessage}. Kluczowa poprawka: pola binarne jako tablice JSON \texttt{[0,1,2,...]}, nie \texttt{json::binary\_t}, które ginęło przy tekstowym \texttt{dump()}.

\subsection{FileTransfer}

\textbf{FileSender} --- statyczne metody wysyłki i \texttt{chunk\_aad()}.\\
\textbf{FileSenderSession} --- stan sesji z buforem chunków do retransmisji.\\
\textbf{FileReceiver} --- mapa aktywnych transferów, logika NACK, \texttt{try\_finalize()}.

Stałe: \texttt{FILE\_CHUNK\_SIZE = 65536}, \texttt{FILE\_MAX\_NACK\_ROUNDS = 5}.

\subsection{ChatApp i ChatAppRender}

\texttt{ChatApp} --- logika: profile, sieć, szyfrowanie wiadomości, pokoje, PFS, transfer plików, benchmarki.\\
\texttt{ChatAppRender.cpp} --- renderowanie ImGui (oddzielenie od logiki).

\textbf{Theme.h} --- autorski motyw kolorystyczny (ciemny, akcent fioletowy).

\subsection{CMake i cele build}

\begin{table}[H]
\centering
\begin{tabular}{@{}llp{7cm}@{}}
\toprule
\textbf{Cel} & \textbf{Typ} & \textbf{Opis} \\
\midrule
\texttt{quill\_core} & STATIC & Krypto + sieć + protokół \\
\texttt{imgui} & STATIC & Dear ImGui + backends (tylko GUI) \\
\texttt{Quill} & EXECUTABLE & Aplikacja (gdy \texttt{QUILL\_BUILD\_GUI=ON}) \\
\texttt{quill\_tests} & EXECUTABLE & Google Test (jeśli znaleziono GTest) \\
\bottomrule
\end{tabular}
\end{table}

Opcja \texttt{QUILL\_BUILD\_GUI=OFF} pomija GLFW, OpenGL i ImGui --- używana w CI.

% ═══════════════════════════════════════════════════════════════════
\section{Interfejs użytkownika}
% ═══════════════════════════════════════════════════════════════════

\subsection{Technologia}

Dear ImGui v1.92.8 (git submodule), backend GLFW + OpenGL 3.3. Renderowanie w pętli głównej \texttt{main.cpp}; dane sieciowe w wątkach roboczych, komunikacja z UI przez kolejki i mutexy.

\subsection{Ekrany i panele}

\begin{table}[H]
\centering
\small
\begin{tabular}{@{}p{3.5cm}p{10cm}@{}}
\toprule
\textbf{Panel} & \textbf{Funkcja} \\
\midrule
Login & Lista profili, tworzenie (passphrase $\times$2), odblokowanie \\
Setup & Tryb Server/Client, host, port, poziom bezpieczeństwa \\
Chat & Wiadomości, wybór pokoju, licznik wiadomości \\
Handshake & Wizualizator kroków z czasami ms, stan TOFU \\
Security Dashboard & Szacunki czasu złamania, rozmiary kluczy \\
Benchmarks & Pomiar FAST/BALANCED/MAX (keygen, encaps, sign...) \\
Tools & Packet Inspector (hex dump), symulator Tamper (MITM) \\
File transfer & Wybór pliku (portable-file-dialogs), postęp, folder pobrań \\
Trusted Peers & Lista known\_hosts, Verify, Remove \\
Profile menu & Fingerprint, logout (zablokowany w sesji) \\
\bottomrule
\end{tabular}
\end{table}

\subsection{Handshake visualizer}

Każdy krok handshake wywołuje callback \texttt{StepCallback(label, time\_ms)}. UI zapisuje listę \texttt{HandshakeStep} i rysuje pasek postępu z czasami, np.:
\begin{itemize}
    \item „Server Signature Verified (ML-DSA)" --- 1.2\,ms,
    \item „KEM Encapsulation (ML-KEM)" --- 0.4\,ms,
    \item „Session Key Derived (HKDF-SHA256)" --- 0.1\,ms.
\end{itemize}

\subsection{Security Dashboard}

Wyświetla szacunki (\texttt{SecurityEstimate}):
\begin{itemize}
    \item czas złamania na superkomputerze klasycznym,
    \item czas złamania na komputerze kwantowym,
    \item werdykt kolorystyczny (bezpieczny / podatny).
\end{itemize}

Dla RSA/ECDH werdykt byłby „podatny"; dla ML-KEM/ML-DSA --- „odporny na kwantowe".

\subsection{Narzędzia audytowe}

\textbf{Packet Inspector} --- podgląd ostatnich pakietów w formie hex + ASCII (edukacyjny, do prezentacji).

\textbf{Tamper simulator} --- celowa modyfikacja bajtu w pakiecie handshake lub sesji, demonstracja odrzucenia przez weryfikację GCM/podpisu.

% ═══════════════════════════════════════════════════════════════════
\section{Scenariusze użycia}
% ═══════════════════════════════════════════════════════════════════

\subsection{Scenariusz 1: Pierwsze uruchomienie (demo lokalne)}

\begin{enumerate}
    \item Uruchom \texttt{./Quill} --- ekran logowania.
    \item Utwórz profil \texttt{Alice} z hasłem (min.\ 8 znaków).
    \item Zaloguj się; zanotuj fingerprint w pasku statusu.
    \item Wybierz tryb \textbf{Server}, port 9999, poziom BALANCED; Start.
    \item W drugiej instancji: profil \texttt{Bob}, tryb \textbf{Client}, localhost:9999.
    \item Handshake automatyczny; TOFU: Bob widzi UNVERIFIED dla serwera.
    \item Wymiana wiadomości w pokoju \texttt{general}.
\end{enumerate}

\subsection{Scenariusz 2: Weryfikacja TOFU}

\begin{enumerate}
    \item Po pierwszym połączeniu porównaj fingerprint serwera out-of-band (np.\ obie instancje na jednym ekranie).
    \item W panelu PEER TRUST kliknij „Mark as Verified".
    \item Przy kolejnym połączeniu stan = VERIFIED.
    \item Symulacja ataku: usuń klucz serwera, wygeneruj nowy profil serwera --- klient dostaje MISMATCH i blokadę.
\end{enumerate}

\subsection{Scenariusz 3: Transfer pliku}

\begin{enumerate}
    \item Po nawiązaniu sesji: przycisk Send File, wybór pliku.
    \item Odbiorca w tym samym pokoju otrzymuje plik w katalogu pobrań.
    \item Postęp: $X/Y$ chunków w UI.
    \item Przy symulowanym dropie chunka: FILE\_NACK i retransmisja (testowane w \texttt{test\_filetransfer.cpp}).
\end{enumerate}

\subsection{Scenariusz 4: Rotacja PFS}

\begin{enumerate}
    \item W trakcie sesji: Rotate Keys (ręcznie) lub poczekaj na auto-rotację serwera.
    \item Nowy handshake na tym samym TCP; liczniki seq zerowane.
    \item Stare wiadomości nie deszyfrują się nowym kluczem.
\end{enumerate}

\subsection{Scenariusz 5: Benchmark poziomów}

\begin{enumerate}
    \item Otwórz panel Benchmarks; uruchom dla FAST, BALANCED, MAX.
    \item Porównaj \texttt{total\_hs\_ms} --- FAST najszybszy, MAX najwolniejszy.
    \item Wyniki do prezentacji: kompromis bezpieczeństwo vs.\ wydajność.
\end{enumerate}

\subsection{Scenariusz 6: Prezentacja na zajęciach}

Typowy przebieg demo ($\sim$15 minut):
\begin{enumerate}
    \item Slajd: motywacja PQC (Shor, NIST 2024).
    \item Uruchomienie dwóch instancji; pokaz fingerprintów.
    \item Handshake visualizer --- omówienie kroków i czasów ms.
    \item Wysłanie wiadomości; Packet Inspector (zaszyfrowany payload).
    \item Tamper simulator --- pokaz odrzucenia przez GCM.
    \item Zmiana klucza serwera --- alert TOFU MISMATCH.
    \item Transfer małego pliku; opcjonalnie Security Dashboard.
\end{enumerate}

% ═══════════════════════════════════════════════════════════════════
\section{Wyniki wydajnościowe (orientacyjne)}
% ═══════════════════════════════════════════════════════════════════

Poniższe wartości pochodzą z wbudowanych benchmarków na typowym laptopie deweloperskim (Linux, Release build). Służą porównaniu poziomów --- nie są gwarancją SLA.

\begin{table}[H]
\centering
\small
\begin{tabular}{@{}lrrrr@{}}
\toprule
\textbf{Poziom} & \textbf{Keygen (ms)} & \textbf{Encaps (ms)} & \textbf{Sign (ms)} & \textbf{Handshake $\Sigma$ (ms)} \\
\midrule
FAST     & $\sim$0.5 & $\sim$0.1 & $\sim$0.5 & $\sim$2--5 \\
BALANCED & $\sim$1.5 & $\sim$0.3 & $\sim$1.5 & $\sim$5--12 \\
MAX      & $\sim$3.0 & $\sim$0.5 & $\sim$2.5 & $\sim$10--25 \\
\bottomrule
\end{tabular}
\caption{Szacunkowe czasy --- zależne od CPU i liboqs. Mierzone w panelu Benchmarks.}
\end{table}

AES-256-GCM na wiadomość tekstową ($<$1\,KB) to zwykle $<$0.1\,ms --- dominują koszty PQC przy handshake, nie szyfrowanie sesji. Dla komunikacji ciągłej po nawiązaniu tunelu opóźnienie jest akceptowalne nawet na poziomie MAX.

\textbf{Wnioski praktyczne:}
\begin{itemize}
    \item \texttt{BALANCED} --- sensowny domyślny wybór (Kyber-768 = rekomendacja NIST),
    \item \texttt{FAST} --- gdy liczy się opóźnienie (np.\ demo na słabszym sprzęcie),
    \item \texttt{MAX} --- gdy priorytetem jest maksymalna siła, kosztem czasu handshake.
\end{itemize}

% ═══════════════════════════════════════════════════════════════════
\section{Testowanie}
% ═══════════════════════════════════════════════════════════════════

\subsection{Strategia testów}

Testy w katalogu \texttt{tests/} linkują wyłącznie \texttt{quill\_core} + Google Test. Brak testów GUI (headless CI). Pokrycie obejmuje:
\begin{itemize}
    \item wektory referencyjne (RFC 5869 HKDF),
    \item testy negatywne (tampering, złe klucze, replay),
    \item testy integracyjne (handshake przez \texttt{socketpair}),
    \item testy wszystkich trzech poziomów bezpieczeństwa (parametryzacja GTest).
\end{itemize}

\subsection{Katalog testów (97 przypadków)}

\subsubsection{test\_hkdf.cpp (6)}
\texttt{Rfc5869TestCase1}, \texttt{Rfc5869TestCase3}, \texttt{ClientServerSymmetry}, \texttt{DomainSeparationByLevel}, \texttt{TranscriptBinding}, \texttt{RejectsEmptyIkm}.

\subsubsection{test\_argon2.cpp (6)}
\texttt{Deterministic}, \texttt{DifferentPassphraseDifferentKey}, \texttt{DifferentSaltDifferentKey}, \texttt{OutputLength}, \texttt{RandomSaltsDiffer}, \texttt{RejectsEmptySalt}.

\subsubsection{test\_aesgcm.cpp (13)}
Roundtrip string/bytes, tamper ciphertext/tag, wrong key, unique nonce, bad key size, AAD roundtrip/wrong/missing/unexpected, AAD bytes, \texttt{ReplayBoundToSeqViaAad}.

\subsubsection{test\_kyber.cpp (10)}
Parametryzacja FAST/BALANCED/MAX: \texttt{EncapsDecapsSharedSecretMatches}, \texttt{FreshKeypairsPerCall}, \texttt{WrongSecretKeyGivesDifferentSecret}; plus \texttt{UnknownLevelThrows}.

\subsubsection{test\_dilithium.cpp (13)}
Parametryzacja $\times$4 testy $\times$3 poziomy + \texttt{LevelAlgorithmMapping}.

\subsubsection{test\_identity.cpp (10)}
Plaintext persist, reload, encrypted save/load, wrong passphrase, encrypted without passphrase, QID1$\rightarrow$QID2 migration, loose permissions, corrupted key, per-algorithm identities, fingerprint format.

\subsubsection{test\_profile.cpp (10)}
Create/unlock, wrong/empty/short/repeated passphrase, duplicate name, path traversal, list profiles, failed create cleanup, unknown profile.

\subsubsection{test\_truststore.cpp (7)}
First contact UNVERIFIED, second KNOWN, mark VERIFIED, MISMATCH preserves entry, remove allows fresh TOFU, distinct peers per host:port, corrupted known\_hosts hard error.

\subsubsection{test\_messageformat.cpp (3)}
Serialize roundtrip, binary payload all byte values, malformed throws.

\subsubsection{test\_filetransfer.cpp (14)}
Multi-chunk roundtrip, corrupted chunk, missing chunk NACK, selective repeat recovery, max NACK exceeded, wrong final hash, duplicate chunks, cross-transfer injection, wrong chunk index, chunk\_hash valid/wrong/missing, make\_NACK indices, SHA3 known vector.

\subsubsection{test\_cryptomanager.cpp (5)}
Same session key both sides, all security levels, re-handshake fresh key + KNOWN trust, server key change blocks, tampered KEM ciphertext rejected.

\subsection{Uruchamianie testów}

\begin{lstlisting}[style=bash]
cd build
cmake -DCMAKE_BUILD_TYPE=Release ..
make -j$(nproc)
./quill_tests                    # bezposrednio
ctest --output-on-failure        # przez CTest
\end{lstlisting}

\subsection{CI GitHub Actions}

Plik \texttt{.github/workflows/ci.yml}:
\begin{enumerate}
    \item checkout z submodułami,
    \item instalacja zależności Ubuntu (\texttt{libssl-dev}, \texttt{libargon2-dev}, \texttt{libgtest-dev}, \texttt{nlohmann-json3-dev}),
    \item budowa i instalacja liboqs ze źródeł,
    \item \texttt{cmake -DQUILL\_BUILD\_GUI=OFF},
    \item \texttt{ctest --output-on-failure}.
\end{enumerate}

Każdy push/PR na \texttt{master}/\texttt{main} weryfikuje warstwę kryptograficzną.

% ═══════════════════════════════════════════════════════════════════
\section{Budowa i wdrożenie}
% ═══════════════════════════════════════════════════════════════════

\subsection{Wymagania systemowe}

\textbf{Kompilator:} C++23 (GCC 13+, Clang 16+).

\textbf{Biblioteki:}
\begin{itemize}
    \item CMake $\geq$ 3.20,
    \item OpenSSL (dowolna wspierana; $<$ 3.2 wymaga libargon2),
    \item liboqs (instalacja ze źródeł),
    \item nlohmann\_json,
    \item Google Test (opcjonalnie, dla testów),
    \item GLFW3 + OpenGL (dla GUI).
\end{itemize}

\subsection{Instalacja liboqs (Arch Linux)}

\begin{lstlisting}[style=bash]
git clone https://github.com/open-quantum-safe/liboqs
cd liboqs && mkdir build && cd build
cmake -DCMAKE_BUILD_TYPE=Release ..
make -j$(nproc)
sudo cmake --install .
\end{lstlisting}

\subsection{Pełna budowa z GUI}

\begin{lstlisting}[style=bash]
git clone --recurse-submodules https://github.com/Smitchelld/Quill.git
cd Quill
mkdir build && cd build
cmake -DCMAKE_BUILD_TYPE=Release ..
make -j$(nproc)
./quill_tests    # 97/97
./Quill          # GUI
\end{lstlisting}

Jeśli brak submodułu ImGui:
\begin{lstlisting}[style=bash]
git submodule update --init --recursive
\end{lstlisting}

\subsection{Pakietowanie (opcjonalne)}

Projekt nie dostarcza oficjalnego pakietu .deb/.pkg. Możliwe kierunki: Docker z \texttt{QUILL\_BUILD\_GUI=OFF} (tylko testy), AppImage z zależnościami GLFW --- poza zakresem v1.0.0.

% ═══════════════════════════════════════════════════════════════════
\section{Analiza bezpieczeństwa --- podsumowanie}
% ═══════════════════════════════════════════════════════════════════

\subsection{Mocne strony}

\begin{itemize}
    \item Zero klasycznej asymetrii w protokole --- zgodność z celem PQC.
    \item Fail-closed: TOFU MISMATCH, replay, GCM tag failure --- wszystko odrzuca, nie degraduje.
    \item Separacja \texttt{quill\_core} --- logika krypto testowalna bez UI.
    \item 97 testów z scenariuszami ataku (MITM, tamper, injection).
    \item Klucze at rest chronione Argon2id (memory-hard).
    \item PFS przez efemeryczny Kyber per sesję/rotację.
\end{itemize}

\subsection{Słabe strony (akceptowane)}

\begin{itemize}
    \item Hub widzi plaintext --- nie E2E.
    \item Brak PKI --- TOFU wymaga dyscypliny użytkownika.
    \item Passphrase jedyna obrona offline --- brak hardware token.
    \item Pliki w RAM przy wysyłce --- limit rozmiaru.
    \item Brak forward secrecy dla tożsamości DSA (trwały klucz) --- standardowe dla podpisów.
\end{itemize}

\subsection{Porównanie z TLS 1.3 + hybryda PQC}

Quill nie implementuje TLS. Różnice:
\begin{itemize}
    \item TLS 1.3: klasyczny ECDHE domyślnie; hybrydy PQC dopiero się pojawiają.
    \item Quill: czysty ML-KEM + ML-DSA od początku; pełna kontrola nad transkryptem HKDF.
    \item TLS: PKI / certyfikaty; Quill: TOFU.
    \item TLS: dojrzały ekosystem; Quill: projekt akademicki, własny protokół.
\end{itemize}

% ═══════════════════════════════════════════════════════════════════
\section{Znane ograniczenia i dług techniczny}
% ═══════════════════════════════════════════════════════════════════

\begin{longtable}{@{}p{2.8cm}p{4cm}p{6.5cm}@{}}
\toprule
\textbf{Obszar} & \textbf{Opis} & \textbf{Wpływ} \\
\midrule
\endhead
E2E vs hub & Serwer odszyfrowuje i re-szyfruje & Brak poufności wobec hosta \\
NAT & Tylko TCP bezpośredni & Brak połączeń przez typowy router domowy bez przekierowania portu \\
PKI & Brak CA & Zaufanie TOFU + out-of-band fingerprint \\
RAM plików & Cały plik w pamięci nadawcy & Limit $\sim$ rozmiar RAM \\
JSON protokół & Brak wersji pola \texttt{version} & Trudniejsza ewolucja bez łamania kompatybilności \\
Wątki UI & Niektóre pola czytane bez mutex & Kosmetyczne race w licznikach UI \\
CI bez GUI & Brak testów renderowania ImGui & Regresje UI tylko manualnie \\
OpenSSL wersje & Dual backend Argon2 & Złożoność build, ale pokryta CI \\
\bottomrule
\caption{Rejestr ograniczeń v1.0.0}
\end{longtable}

% ═══════════════════════════════════════════════════════════════════
\section{Podsumowanie i dalsza praca}
% ═══════════════════════════════════════════════════════════════════

\subsection{Podział prac (przykładowy)}

W projekcie zespołowym (2 osoby) naturalny podział:
\begin{itemize}
    \item \textbf{Warstwa krypto + protokół:} CryptoManager, TOFU, HKDF, testy integracyjne,
    \item \textbf{Sieć + UI:} NetworkServer/Client, ChatApp, ImGui, transfer plików,
    \item \textbf{Wspólne:} review bezpieczeństwa, dokumentacja, prezentacja.
\end{itemize}

\noindent\textit{[Uzupełnij rzeczywistym podziałem w zespole.]}

\subsection{Osiągnięcia}

Quill v1.0.0 dostarcza działający komunikator post-kwantowy z:
\begin{itemize}
    \item pełnym handshake ML-KEM + ML-DSA/Falcon,
    \item trzema poziomami bezpieczeństwa,
    \item szyfrowaną sesją AES-256-GCM z anty-replay,
    \item TOFU fail-closed,
    \item profilami i szyfrowaniem kluczy Argon2id,
    \item transferem plików z SHA-3 i selective repeat,
    \item interfejsem ImGui z narzędziami edukacyjnymi,
    \item 97 testami i CI.
\end{itemize}

Projekt jest gotowy do prezentacji i oddania w ramach Studio Projektowego, z jawnie opisanymi ograniczeniami architektonicznymi.

\subsection{Planowana praca inżynierska}

Temat: \emph{inteligentny dobór poziomu PQC} (FAST/BALANCED/MAX) w zależności od:
\begin{itemize}
    \item opóźnienia sieci (RTT),
    \item mocy obliczeniowej urządzenia,
    \item wymaganego poziomu bezpieczeństwa (klasyfikacja danych),
    \item zużycia baterii (urządzenia mobilne/IoT).
\end{itemize}

Quill dostarcza już benchmarki per poziom --- naturalny punkt wyjścia do heurystyk lub modelu ML.

\subsection{Roadmapa techniczna}

\begin{enumerate}
    \item NAT traversal (UDP hole punching + rendezvous) --- jedyny duży brakujący feature,
    \item opcjonalnie: ASan/UBSan w CI,
    \item opcjonalnie: streaming plików bez pełnego buforu RAM,
    \item opcjonalnie: wersjonowanie protokołu JSON,
    \item długoterminowo: prawdziwy E2E (double ratchet / per-recipient keys).
\end{enumerate}

% ═══════════════════════════════════════════════════════════════════
\appendix
\section{Załącznik A: Zmienne środowiskowe}
% ═══════════════════════════════════════════════════════════════════

\begin{table}[H]
\centering
\begin{tabular}{@{}ll@{}}
\toprule
\textbf{Zmienna} & \textbf{Domyślnie} \\
\midrule
\texttt{QUILL\_PROFILES\_DIR} & \texttt{\textasciitilde/.quill/profiles} \\
\texttt{QUILL\_IDENTITY\_DIR} & \texttt{\textasciitilde/.quill/identity} \\
\bottomrule
\end{tabular}
\end{table}

\section{Załącznik B: Rozmiary typowych artefaktów kryptograficznych}

\begin{table}[H]
\centering
\small
\begin{tabular}{@{}llr@{}}
\toprule
\textbf{Poziom} & \textbf{Obiekt} & \textbf{Ok.\ rozmiar (B)} \\
\midrule
BALANCED & Kyber-768 public key & $\sim$1184 \\
BALANCED & Kyber-768 ciphertext & $\sim$1088 \\
BALANCED & ML-DSA-65 signature & $\sim$3309 \\
BALANCED & ML-DSA-65 public key & $\sim$1952 \\
\midrule
Sesja & AES-256 key & 32 \\
Sesja & GCM nonce & 12 \\
Sesja & GCM tag & 16 \\
\bottomrule
\end{tabular}
\caption{Wartości orientacyjne --- dokładne rozmiary zależą od liboqs.}
\end{table}

\section{Załącznik C: Konwencje nazewnictwa peer\_id w TOFU}

\begin{table}[H]
\centering
\begin{tabular}{@{}ll@{}}
\toprule
\textbf{Kierunek} & \textbf{Format peer\_id} \\
\midrule
Klient identyfikuje serwer & \texttt{srv:localhost:9999} \\
Serwer identyfikuje klienta & \texttt{cli:<SHA3-fingerprint>} \\
\bottomrule
\end{tabular}
\end{table}

\section{Załącznik D: Przykład zaszyfrowanej wiadomości CHAT (logiczny)}

Po szyfrowaniu wiadomość na linii ma postać JSON (wartości przykładowe):

\begin{lstlisting}[style=json]
{
  "type": "CHAT",
  "sender": "Alice",
  "seq": 42,
  "timestamp": 1718123456789,
  "nonce": [/* 12 losowych bajtow */],
  "payload": [/* ciphertext + tag GCM, binarnie jako tablica */]
}
\end{lstlisting}

Przy deszyfracji odbiorca:
\begin{enumerate}
    \item buduje AAD = \texttt{"CHAT|42"},
    \item wywołuje \texttt{AesGcm::decrypt(key, nonce, payload, aad)},
    \item sprawdza \texttt{seq > recv\_seq}, aktualizuje \texttt{recv\_seq}.
\end{enumerate}

\section{Załącznik E: Format pliku known\_hosts}

\begin{lstlisting}[style=json]
{
  "peers": [
    {
      "peer_id": "srv:localhost:9999",
      "fingerprint": "a3f29c1b44de...",
      "algo": "ML-DSA-65",
      "verified": true
    }
  ]
}
\end{lstlisting}

Plik zapisywany atomowo, chmod 0600. Uszkodzony JSON powoduje twardy błąd --- brak cichego nadpisania (test \texttt{CorruptedKnownHostsIsHardError}).

\section{Załącznik F: Lista plików źródłowych quill\_core}

\begin{lstlisting}[style=bash]
src/crypto/KyberKEM.cpp
src/crypto/DilithiumSign.cpp
src/crypto/AesGcm.cpp
src/crypto/Hkdf.cpp
src/crypto/Argon2.cpp
src/crypto/IdentityManager.cpp
src/crypto/ProfileManager.cpp
src/crypto/TrustStore.cpp
src/crypto/CryptoManager.cpp
src/network/Socket.cpp
src/network/NetworkServer.cpp
src/network/NetworkClient.cpp
src/protocol/MessageFormat.cpp
src/protocol/FileTransfer.cpp
\end{lstlisting}

% ═══════════════════════════════════════════════════════════════════
\begin{thebibliography}{99}

\bibitem{fips203}
NIST FIPS 203 --- Module-Lattice-Based Key-Encapsulation Mechanism Standard.
\url{https://csrc.nist.gov/pubs/fips/203/final}

\bibitem{fips204}
NIST FIPS 204 --- Module-Lattice-Based Digital Signature Standard.
\url{https://csrc.nist.gov/pubs/fips/204/final}

\bibitem{fips205}
NIST FIPS 205 --- Stateless Hash-Based Digital Signature Standard (SPHINCS+).
\url{https://csrc.nist.gov/pubs/fips/205/final}

\bibitem{liboqs}
Open Quantum Safe Project --- liboqs documentation.
\url{https://openquantumsafe.org}

\bibitem{kyber}
J. Bos, L. Ducas, E. Kiltz, T. Lepoint, V. Lyubashevsky, J. Schanck, P. Schwabe, G. Seiler, D. Stehlé.
\emph{CRYSTALS-Kyber: A CCA-Secure Module-Lattice-Based KEM.}

\bibitem{dilithium}
L. Ducas, E. Kiltz, T. Lepoint, V. Lyubashevsky, P. Schwabe, G. Seiler, D. Stehlé.
\emph{CRYSTALS-Dilithium: A Lattice-Based Digital Signature Scheme.}

\bibitem{rfc5869}
H. Krawczyk. RFC 5869 --- HMAC-based Extract-and-Expand Key Derivation Function (HKDF). 2010.

\bibitem{rfc9106}
A. Biryukov, D. Dinu, D. Khovratovich, S. Josefsson.
RFC 9106 --- Argon2 Memory-Hard Function for Password Hashing and Proof-of-Work. 2021.

\bibitem{ssh-tofu}
T. Ylonen, C. Lonvick. RFC 4251 --- The Secure Shell (SSH) Protocol Architecture. 2006.

\bibitem{nist-pqc}
NIST Post-Quantum Cryptography Standardization.
\url{https://csrc.nist.gov/projects/post-quantum-cryptography}

\end{thebibliography}

\end{document}
